% Diese Fan-Geschichte ist noch in Arbeit. Ich teile sie bereits hier, damit sie für mich nicht verloren geht, falls mein Laptop spontan verschwinden sollte, und damit sie für euch nicht völlig verloren geht, falls ich spontan verschwinden sollte. Fühlt euch in letzterem Fall vollkommen frei, hierrin zu schmökern und nach Lust und Laune weiter daran zu werkeln. Aber solange ich ein aktiver Andori bleibe, lest bitte nicht zu genau nach, was ich geplant haben könnte. LG BBB :D

\begin{chapterbox}
    \chapter{Eine Fährtenleserin entscheidet sich (WIP)}
    \label{wip-eine-faehrtenleserin-entscheidet-sich}
    \az{Jahr 75}
    
    \begin{center}
        Teil D der Neuen Abenteuer\sidenotemark

        Fortsetzung von \hyperref[wip-in-vertrauten-landen]{in vertrauten Landen (WIP)}
    \end{center}

    Fenn schmort im Kerker der Rietburg vor sich hin. Als die Armee der Krahder sich der Burg nähert, ergibt sich eine letzte Möglichkeit, seine Fähigkeiten für seine Befreiung zu nutzen.
\end{chapterbox}


\sidenotetext{Anders als bei den vorherigen Neuen Abenteuern steht hier nicht explizit, von \gend{welcher Held*in}{welchem Helden} dieses Abenteuer handelt. Dies ist so, weil \siehechronik{Ein Rabe entscheidet sich (2019)} die einzige Geschichte von Peter Gustav Bartschat ist, die auf dem Titelbild nicht explizit schreibt, von \gend{welcher Held*in}{welchem Helden} sie handelt. Dies ist der einzige Grund. Es gibt keinen anderen Grund. Oder etwa doch?}


\section{Prolog}
\az{Jahr 74}

Fenn suchte Arbons Eltern auf.

Fenn wurde vor den Rat der \gend{Bewahrer*innen}{Bewahrer} gestellt.

Ken Dorr, der neue Regent von Andor, hatte Interesse an seinen Fähigkeiten und verlangte seine Auslieferung.




\newpage
\section{Kapitel}
\az{Jahr 75}

Ken besuchte Fenn oft im Kerker.






\begin{commentbox}
    Weiter geht es in \hyperref[wip-die-wahre-macht-des-zweiten-schilds]{Die wahre Macht des zweiten Schilds (WIP)}.
\end{commentbox}